\documentclass[a4paper, oneside]{article}
\usepackage{graphicx}
\usepackage{amsthm}
\usepackage{amsmath}
\usepackage{amssymb}
\usepackage[a4paper,
            bindingoffset=0.2in,
            left=2cm,
            right=2cm,
            top=2cm,
            bottom=2cm,
            footskip=.25in]{geometry}
\usepackage[italian]{babel}
\usepackage{pgfplots}
\usepackage{tabularx}
\usepackage{tikz}
\usepackage{wrapfig}
\usepackage{color}
\usepackage[d]{esvect}
\definecolor{page}{rgb}{0.129,0.157,0.212}
\pagecolor{page}
\color{white}
\graphicspath{ {./images/} }
\usetikzlibrary{shapes.geometric}
\usetikzlibrary{datavisualization}
\usetikzlibrary{datavisualization.formats.functions}
\usetikzlibrary{patterns}
\pgfplotsset{width=10cm,compat=1.18}

\title{Appunti di Lab II}
\author{Tommaso Miliani}
\date{16-03-26}

\begin{document}
\newtheoremstyle{theoremEnv}
                {}          % Space above
                {}          % Space below
                {\slshape}  % Body font
                {}          % Indent amount
                {\bfseries} % Head font
                {.}         % Punctuation after head
                {\newline}  % Space after theorem head
                {}          % Theorem head spec
\theoremstyle{theoremEnv}

\newtheorem{definition}{Definizione}[section]
\newtheorem{theorem}{Teorema}[section]
\newtheorem{lemma}{Proposizione}[section]
\newtheorem{observation}{Osservazione}[section]
\newtheorem{corollary}{Corollario}[theorem]
\newtheorem{example}{Esempio}[section]
\newtheorem{remark}{Enunciato}[section]

\maketitle

\section{Fine esperienza 1 perché si}
Si trova anche tutto su moodle 


\section{Esperienza 2}
Si utilizza un amplificatore che è costituito come in figura.
Lo scopo di questa esperienza è la determinazione del guadagno 
in continua e la caratterizzazione delle caratteristiche dell'amplificatore. 
Gli amplificatori come questi permettono di amplificare il segnale in 
ingresso, per cui all'uscita si osserva che
\begin{gather*}
    V_0 = A(V_+ - V_-)
\end{gather*}
L'amplificatore in laboratorio è un LF156, per cui i valori tipici sono 
\begin{itemize}
    \item $A = 10^{5}\sim 10^{6}$: è molto variabile e varia sia con
    la temperatura che da chip a chip
    \item $R_i = 10^{11} \Omega$: resistenza in ingresso
    \item $R_{out} = 200\Omega$: è la resistenza in uscita. 
\end{itemize}
Per poter stabilizzare questo oggetto (è instabile perché dipende a cosa è 
attaccato) si collega l'uscita in retroazione all'ingresso $-$. Nell'esperienza
si deve dunque realizzare una configurazione ad amplificatore operazionale in modo 
tale che parte del segnale in uscita sia rimandato in ingresso.

\subsection{Parte I: stabilizzare l'amplificatore}
L'amplificazione a circuito-retroazionato è data solo e soltanto 
(nel limite in cui $R_f$ e $R_s$ siano molto minori di $A$):
\begin{gather*}
    A_f = - \frac{R_f}{R_s}
\end{gather*}
IN questo modo si realizza un sistema molto più stabile (il segno negativo è 
perché il segnale rientra dal polo negativo) e dipende dal valore delle resistenze. 
Questo oggetto permette di ripetere e amplificare questo segnale in modo tale che
la resistenza in uscita tenda a zero.  \\
Si utilizza una breadboard (le cui linee sono collegate tra loro)
e quelle centrali hanno le colonne collegate tra loro. Per fare 



\end{document}